\chapter{Hỏi Khi Lớp Đang Ngủ}
\chapterintro{Câu hỏi phải đủ ngắn để kéo mắt lên}{Lúc lớp mệt, câu hỏi càng ngắn và càng gần thì càng dễ cứu nhịp.}

\section{Nếu bây giờ em phải nói đúng 3 từ về trạng thái của mình, em nói gì?}
\begin{cauhoibox}{Câu hỏi}
Nếu bây giờ em phải nói đúng 3 từ về trạng thái của mình, em nói gì?
\end{cauhoibox}
\begin{taisaocoich}
Câu hỏi ngắn, dễ trả lời, ít áp lực. Nó giúp lớp bật nói mà không cần động não quá nặng ngay lúc đang mệt.
\end{taisaocoich}
\begin{goiybien}
Hợp đầu tiết sáng, đầu tiết chiều, hoặc ngay sau giờ ra chơi.
\end{goiybien}
\ngancach

\section{Nếu cho tiết này một nút bấm để tỉnh lại, em bấm gì?}
\begin{cauhoibox}{Câu hỏi}
Nếu cho tiết này một nút bấm để tỉnh lại, em bấm gì?
\end{cauhoibox}
\begin{taisaocoich}
Câu hỏi có hình ảnh nên học sinh thích trả lời. Nó cũng giúp thầy cô đo nhanh lớp đang cần đổi nhịp kiểu nào.
\end{taisaocoich}
\begin{goiybien}
Sau khi nghe 2 đến 3 ý, chọn luôn một cách đổi nhịp thật ngắn cho lớp.
\end{goiybien}
\ngancach

\section{Điều gì khiến em thức dậy trong một tiết học?}
\begin{cauhoibox}{Câu hỏi}
Điều gì khiến em thức dậy trong một tiết học: câu hỏi, ví dụ vui, hay hoạt động ngắn?
\end{cauhoibox}
\begin{taisaocoich}
Đây là câu hỏi vừa kéo lớp nói, vừa cho thầy cô dữ liệu thật về cách lớp phản ứng với bài giảng.
\end{taisaocoich}
\begin{goiybien}
Rất hợp dùng đầu năm hoặc khi muốn điều chỉnh cách dạy cho một lớp cụ thể.
\end{goiybien}
\ngancach

\section{Nếu lớp được đổi tư thế một phút, em muốn đổi kiểu nào?}
\begin{cauhoibox}{Câu hỏi}
Nếu lớp được đổi tư thế một phút để tỉnh lại, em muốn đổi kiểu nào?
\end{cauhoibox}
\begin{taisaocoich}
Lúc lớp mệt, câu hỏi gắn với cơ thể thường kéo phản hồi nhanh hơn câu hỏi nặng về nội dung. Nó cũng cho thầy cô gợi ý để đổi nhịp thật ngắn.
\end{taisaocoich}
\begin{goiybien}
Sau khi nghe 2 ý, chọn ngay một cách đơn giản như đứng lên, đổi chỗ tay, hay quay lại nói với bạn bên cạnh một câu.
\end{goiybien}
\ngancach

\section{Nếu cơn buồn ngủ biết nói, nó đang nói gì với tiết này?}
\begin{cauhoibox}{Câu hỏi}
Nếu cơn buồn ngủ biết nói, nó đang nói gì với tiết này?
\end{cauhoibox}
\begin{taisaocoich}
Câu hỏi có chút nhân hóa nên lớp dễ bật cười và bật nói. Nó giúp thừa nhận sự mệt mà không làm không khí nặng thêm.
\end{taisaocoich}
\begin{goiybien}
Hợp cuối buổi hoặc đầu giờ chiều. Nghe 2 đến 3 câu trả lời rồi chuyển ngay sang một hoạt động ngắn để kéo lớp tỉnh lên.
\end{goiybien}
